\documentclass[12pt]{article}
\usepackage[utf8]{inputenc}
\usepackage[T1]{fontenc}
\usepackage[margin=1in]{geometry}
\usepackage{setspace}
\usepackage{xcolor}
\usepackage[colorlinks=true, urlcolor=blue!60!black, linkcolor=blue!60!black]{hyperref}
\setlength{\parindent}{0em}
\setlength{\parskip}{0.45em}
\renewcommand{\baselinestretch}{1.0}
\begin{document}
\begin{center}
{\Large \textit{Everson v. Board of Education}, 330 U.S. 1 (1947)\par}
\vspace{0.3em}
{\normalsize Case Brief --- One-Page Edition\par}
\end{center}

\textbf{Memory Jogger.} The parochial-school bus case: the Establishment Clause is incorporated, but neutral welfare benefits survive (\href{https://supreme.justia.com/cases/federal/us/330/1/}{330 U.S. 1}).

\textbf{Facts.} Ewing Township, New Jersey, reimbursed parents from tax funds for public-bus fares to public high schools and four Catholic parochial schools, which met state educational requirements but also gave regular religious instruction. Taxpayer Arch R. Everson sued.

\textbf{Procedural History.} New Jersey's highest court upheld the program; Everson appealed.

\textbf{Votes.} Affirmed, 5--4. Black wrote for the Court; Jackson and Rutledge dissented.

\textbf{Holding.} The Establishment Clause binds the states, but a neutral benefit paid to parents of all faiths, like police and fire protection, does not breach the wall between church and state.

\textbf{Reasoning.} Drawing on Madison and Jefferson, Black declared that ``No tax in any amount, large or small, can be levied to support any religious activities or institutions,'' yet held that the state must be neutral, not religion's adversary, and cannot exclude believers from public welfare benefits. Jackson found the rhetoric discordant with the result; Rutledge saw no line between buses and books or salaries.

\textbf{Examples.} \emph{Same-Side:} free transit passes for all students. \emph{Opposite-Side:} tax grants to all faiths for clergy salaries, which fund the religious function itself. \emph{Fence-Sitter:} nonsectarian requirements for charter schools.

\textbf{Critique.} Critics question incorporation and the reliance on Virginia history; Jackson's warning that government regulates what it subsidizes still frames aid disputes.

\textbf{Key Quotation.} ``The `establishment of religion' clause of the First Amendment means at least this: neither a state nor the Federal Government can set up a church. Neither can pass laws which aid one religion, aid all religions, or prefer one religion over another.'' --- Justice Black, Opinion of the Court, 330 U.S. at 15--16.

\end{document}
